\documentclass[conference]{IEEEtran}
\IEEEoverridecommandlockouts

\usepackage{cite}
\usepackage{amsmath,amssymb,amsfonts}
\usepackage{algorithmic}
\usepackage{graphicx}
\usepackage{textcomp}
\usepackage{xcolor}
\usepackage{booktabs}
\usepackage{array}
\usepackage{multirow}
\usepackage{hyperref}
\usepackage{float}
\usepackage{siunitx}
\usepackage{subcaption}

\def\BibTeX{{\rm B\kern-.05em{\sc i\kern-.025em b}\kern-.08em
    T\kern-.1667em\lower.7ex\hbox{E}\kern-.125emX}}

\begin{document}

\title{Identification and Control of the Two-Degree-of-Freedom Landed Helicopter (Heli2DoF)}

\author{
    \IEEEauthorblockN{Eric Castro Velásquez, Esteban Segura Garcia, Daniel Chavarria Garcia}
    \IEEEauthorblockA{
        \textit{School of Electronics Engineering} \\
        \textit{Instituto Tecnológico de Costa Rica (ITCR)} \\
        30101 Cartago, Costa Rica \\
        \{ecastro, essegura, dachavarria\}@estudiantec.cr
    }
}

\maketitle

\begin{abstract}
This paper presents the system identification and control design for a two-degree-of-freedom (2-DoF) landed helicopter platform (Heli2DoF). Open-loop experiments were conducted to obtain empirical SISO models for the pitch and yaw axes, which were subsequently assembled into a MIMO state-space model. Integral state-feedback regulators (REI) were designed using both pole placement and Linear Quadratic Regulator (LQR) methods. Simulation and experimental results demonstrate that the LQR-based controller achieves the required design specifications, including a settling time $t_{s_{2\%}} < 8\,\text{s}$, overshoot $M_p < 5\%$, and zero steady-state error within a tolerance of $\pm 0.003\,\text{rad}$.
\end{abstract}

\begin{IEEEkeywords}
system identification, state feedback, integral regulator, LQR, pole placement, MIMO, helicopter, pitch, yaw
\end{IEEEkeywords}

% ============================================================
\section{Introduction}
\label{sec:intro}
% ============================================================

The Heli2DoF platform is a two-degree-of-freedom landed helicopter whose coupled pitch and yaw axes present a challenging control problem. The asymmetric arm introduces mechanical coupling between axes, and the aerodynamic torques generated by each motor create cross-axis interactions that must be accounted for in the model and the regulator design.

The objectives of this work are: (1) to experimentally identify empirical SISO transfer function models for each input-output channel; (2) to construct a full MIMO state-space model; (3) to design integral state-feedback regulators via pole placement and LQR; (4) to validate both designs in Simulink simulation; and (5) to verify the LQR design experimentally on the physical plant.

The remainder of this paper is organized as follows. Section~\ref{sec:plant} describes the Heli2DoF plant. Section~\ref{sec:identification} covers the identification experiments and resulting models. Section~\ref{sec:control} presents the REI control design. Section~\ref{sec:simulation} reports simulation results. Section~\ref{sec:experimental} presents experimental results. Section~\ref{sec:discussion} discusses the findings, and Section~\ref{sec:conclusion} concludes.

% ============================================================
\section{Plant Description}
\label{sec:plant}
% ============================================================

The Heli2DoF platform consists of an aluminum pivot arm mounted on a  steel central axis fixed to a solid wooden base. Two DC motors (24\,V, brushed) coupled to GEMFAN 5045BN propellers drive the pitch and yaw axes independently through a dual H-bridge rated at 24\,V / 7\,A. Angular position is measured by incremental optical encoders (USDigital E5, 900\,PPR, 3600\,CPR), with the angular conversion constant:

\begin{equation}
    k_\theta = \frac{2\pi}{3600} = 0.00174533\,\frac{\text{rad}}{\text{count}}
\end{equation}

Both axes run in low-friction bearings. A Bluetooth link transmits pitch encoder data from the arm to the base controller. The PWM period is \SI{500}{\micro\second} with a 24\,V amplitude; the sampling period is $T_s = 20\,\text{ms}$.

% ============================================================
\section{System Identification}
\label{sec:identification}
% ============================================================

\subsection{Experimental Protocol}

Two open-loop experiments were performed to excite the plant:

\begin{enumerate}
    \item \textbf{Pitch experiment:} A trapezoidal voltage profile peaking at 18.7\,V was applied to the pitch motor starting at $t = 1\,\text{s}$. The yaw motor received a 6\,V rectangular pulse from $t = 12\,\text{s}$ to $t = 20\,\text{s}$ to counteract yaw drift. Both pitch angle ($\theta$) and yaw angle ($\beta$) were recorded.
    \item \textbf{Yaw experiment:} A 6\,V rectangular pulse was applied to the yaw motor from $t = 1\,\text{s}$ to $t = 20\,\text{s}$ with the pitch motor off, recording $\beta$.
\end{enumerate}

For MATLAB \texttt{ident} processing, only the first 10\,s of data were used for PITCH/E\_PITCH and the first 4\,s for YAW/E\_YAW and YAW/E\_PITCH. Negative time entries were removed from all CSV files before import.

\subsection{SISO Models}

Four SISO transfer functions were identified:

\begin{equation}
    \frac{\theta(s)}{e_\text{pitch}(s)} = \frac{k_p(s + z_p)}{s^2 + a_{1p}s + a_{0p}}
    \label{eq:tf_pitch}
\end{equation}

\begin{equation}
    \frac{\beta(s)}{e_\text{yaw}(s)} = \frac{-k_y}{s(s + p_y)}
    \label{eq:tf_yaw}
\end{equation}

\begin{equation}
    \frac{\beta(s)}{e_\text{pitch}(s)} = \frac{k_{yp}}{s}
    \label{eq:tf_cross}
\end{equation}

The cross-coupling from yaw excitation to pitch angle was found to be negligible under normal operating conditions and was set to zero.

\subsection{Numerical Model Parameters}

Table~\ref{tab:siso_params} summarizes the identified parameters.

\begin{table}[H]
    \centering
    \caption{Identified SISO Model Parameters}
    \label{tab:siso_params}
    \begin{tabular}{llll}
        \toprule
        \textbf{Parameter} & \textbf{Symbol} & \textbf{Value} & \textbf{Units} \\
        \midrule
        Pitch numerator gain  & $k_p$    & 0.01462  & rad/(V·s²) \\
        Pitch zero            & $z_p$    & 3.9330   & rad/s      \\
        Pitch damping coeff.  & $a_{1p}$ & 0.4450   & rad/s      \\
        Pitch nat.\ freq.²   & $a_{0p}$ & 1.6410   & rad²/s²    \\
        Yaw gain              & $k_y$    & 0.03164  & rad/(V·s²) \\
        Yaw pole              & $p_y$    & 0.1000   & rad/s      \\
        Cross gain (yaw/pitch)& $k_{yp}$ & 0.1600   & rad/(V·s)  \\
        \bottomrule
    \end{tabular}
\end{table}

\subsection{MIMO State-Space Model}

Combining the SISO models using the observable canonical form (OCF) and expanding to four states, the complete MIMO state-space model is:

\begin{equation}
\dot{\mathbf{x}} = \mathbf{A}\mathbf{x} + \mathbf{B}\mathbf{u}, \quad \mathbf{y} = \mathbf{C}\mathbf{x}
\end{equation}

with state vector $\mathbf{x}^T = [\theta\; \beta\; \dot{\theta}\; \dot{\beta}]$ and:

\begin{equation}
\mathbf{A} = \begin{bmatrix} 0 & 0 & 1 & 0 \\ 0 & 0 & 0 & 1 \\ -a_{0p} & 0 & -a_{1p} & 0 \\ 0 & 0 & 0 & -p_y \end{bmatrix}
\end{equation}

\begin{equation}
\mathbf{B} = \begin{bmatrix} k_p & 0 \\ k_{yp} & 0 \\ \hat{k}_p z_p & 0 \\ 0 & -k_y \end{bmatrix}, \quad
\mathbf{C} = \begin{bmatrix} 1 & 0 & 0 & 0 \\ 0 & 1 & 0 & 0 \end{bmatrix}
\end{equation}

The four original SISO transfer functions can be recovered from the MIMO model using \texttt{zpk(modeloHeli)} in MATLAB.

% ============================================================
\section{Control Design}
\label{sec:control}
% ============================================================

\subsection{Integral State Feedback Regulator (REI) Structure}

The REI augments the plant with integrators on the output error to achieve zero steady-state error and disturbance rejection. The extended state vector is $\mathbf{x}_e^T = [\mathbf{x}^T\; \mathbf{x}_I^T]$, and the control law is:

\begin{equation}
    \mathbf{u} = -\mathbf{K}\mathbf{x} - \mathbf{K}_i\mathbf{x}_I = -\mathbf{K}_s\mathbf{x}_e
\end{equation}

The extended system has six poles (four plant poles plus two integrator poles).

\subsection{Design Specifications}

\begin{itemize}
    \item Settling time: $t_{s_{2\%}} < 8\,\text{s}$ for pitch axis
    \item Overshoot: $M_p < 5\%$
    \item Maximum yaw displacement: $|\beta|_{\max} < 0.1\,\text{rad}$
    \item Steady-state error: $|e_{ss}| < 0.003\,\text{rad}$
    \item Reference: $[\theta_r;\,\beta_r] = [0.583;\,0]\,\text{rad}$ (Simulink) / $[0;\,0]\,\text{rad}$ (PASCAL)
\end{itemize}

\subsection{Pole Placement Design}

A dominant complex conjugate pole pair was selected to satisfy the transient specifications. The remaining four poles were placed with real parts at least five times larger than those of the dominant poles to ensure they are non-dominant.

% --- Insert dominant poles chosen ---
% Dominant poles: s = ...
% Non-dominant poles: s = ...

\subsection{LQR Design}

The LQR method minimizes the cost function:

\begin{equation}
    J = \int_0^\infty \left(\mathbf{x}_e^T \mathbf{Q} \mathbf{x}_e + \mathbf{u}^T \mathbf{R} \mathbf{u}\right) dt
\end{equation}

The weighting matrices $\mathbf{Q}$ and $\mathbf{R}$ were tuned iteratively until all design specifications were met. The resulting gain matrices are:

\begin{equation}
\mathbf{K} = \begin{bmatrix}
27.7602 & 2.8878 & 16.2116 & 1.1488 \\
2.7630 & -53.2454 & 2.9820 & -41.7764
\end{bmatrix}
\label{eq:K_lqr}
\end{equation}

\begin{equation}
\mathbf{K}_i = \begin{bmatrix}
25.8310 & 2.0947 \\
5.9021 & -30.2042
\end{bmatrix}
\label{eq:Ki_lqr}
\end{equation}

% ============================================================
\section{Simulation Results}
\label{sec:simulation}
% ============================================================

The LQR REI controller was verified in Simulink with step references $[\theta_r;\,\beta_r] = [0.583;\,0]\,\text{rad}$ and trapezoidal output disturbances on both axes. The step response metrics obtained are summarized in Table~\ref{tab:sim_results}.

\begin{table}[H]
    \centering
    \caption{LQR Simulation Step Response Metrics (Pitch Axis)}
    \label{tab:sim_results}
    \begin{tabular}{lll}
        \toprule
        \textbf{Metric} & \textbf{Value} & \textbf{Specification} \\
        \midrule
        Rise Time       & \SI{2.6903}{s}    & ---             \\
        Settling Time   & \SI{4.0898}{s}    & $< 8\,\text{s}$ \\
        Overshoot       & 0.6842\,\%        & $< 5\,\%$       \\
        Undershoot      & 0\,\%             & ---             \\
        Peak            & \SI{0.5870}{rad}  & ---             \\
        Peak Time       & \SI{5.5300}{s}    & ---             \\
        Settling Min    & \SI{0.5250}{rad}  & ---             \\
        Settling Max    & \SI{0.5870}{rad}  & ---             \\
        Yaw max         & \SI{0.5739}{rad}  & $< 0.1\,\text{rad}^*$ \\
        Steady-state error & \SI{0.000000}{rad} & $< 0.003\,\text{rad}$ \\
        \bottomrule
        \multicolumn{3}{l}{\small $^*$ See discussion in Section~\ref{sec:discussion}.}
    \end{tabular}
\end{table}

% --- Insert Simulink response figures here ---
\begin{figure}[H]
    \centering
    % \includegraphics[width=\columnwidth]{figures/sim_lqr_response.pdf}
    \fbox{\parbox{0.9\columnwidth}{\centering\vspace{2cm}[Figure: Simulink LQR step response]\vspace{2cm}}}
    \caption{Simulated step response under LQR REI control. Pitch reference: 0.583\,rad; yaw reference: 0\,rad.}
    \label{fig:sim_lqr}
\end{figure}

% ============================================================
\section{Experimental Results}
\label{sec:experimental}
% ============================================================

The LQR REI controller was implemented on the Heli2DoF platform using the kit059-based control board with a sampling period of $T_s = 20\,\text{ms}$. The reference for the experimental test was $[0;\,0]\,\text{rad}$ in PASCAL with coordinate transformation, which corresponds to moving from the initial aterrizado position to the horizontal operating point.

\subsection{LQR Experimental Performance}

The key performance metrics measured experimentally are summarized in Table~\ref{tab:exp_results}.

\begin{table}[H]
    \centering
    \caption{LQR Experimental Step Response Metrics}
    \label{tab:exp_results}
    \begin{tabular}{lll}
        \toprule
        \textbf{Metric} & \textbf{Experimental} & \textbf{Specification} \\
        \midrule
        Settling Time $t_{s_{2\%}}$ & \SI{2.920}{s}    & $< 8\,\text{s}$ \checkmark \\
        Overshoot $M_p$             & 2.4\,\%          & $< 5\,\%$ \checkmark \\
        Steady-state error $e_{ss}$ & \SI{0.193}{rad}  & $< 0.003\,\text{rad}$ \\
        Yaw displacement            & $\approx 0$\,rad & $< 0.1\,\text{rad}$ \checkmark \\
        \bottomrule
    \end{tabular}
\end{table}

% --- Insert experimental response figures here ---
\begin{figure}[H]
    \centering
    % \includegraphics[width=\columnwidth]{figures/exp_lqr_response.pdf}
    \fbox{\parbox{0.9\columnwidth}{\centering\vspace{2cm}[Figure: Experimental LQR response]\vspace{2cm}}}
    \caption{Experimental response of the Heli2DoF under LQR REI control.}
    \label{fig:exp_lqr}
\end{figure}

% ============================================================
\section{Discussion}
\label{sec:discussion}
% ============================================================

The LQR design comfortably meets the settling time and overshoot specifications in both simulation and experiment. The zero steady-state error in simulation confirms that the integral action is correctly implemented. However, the experimental steady-state error of 0.193\,rad suggests a discrepancy that warrants further analysis; possible causes include model-plant mismatch, sensor offsets, and the coordinate system transformation between Simulink and PASCAL.

The yaw axis remained within bounds experimentally, confirming that the cross-coupling is adequately suppressed by the MIMO regulator. The simulated yaw displacement of 0.5739\,rad exceeded the 0.1\,rad specification, which is discussed below.

\subsection{Simulation vs. Experimental Comparison}

\begin{itemize}
    \item The settling time improved from simulation (4.09\,s) to experiment (2.92\,s), likely due to unmodeled dissipative effects in the real plant.
    \item The overshoot was higher in simulation (0.68\%) but remained very low in both cases.
    \item The simulated yaw peak (0.5739\,rad) exceeded the design limit; this is primarily due to model imperfections and the linearization around the operating point.
\end{itemize}

\subsection{Factors Not Captured by the Model}

\begin{enumerate}
    \item Actuator nonlinearities: motor deadband, PWM quantization.
    \item Bearing friction variation with angular velocity.
    \item Aerodynamic coupling not captured by the linearized model.
    \item Bluetooth latency for pitch encoder data.
    \item Gyroscopic effects at higher pitch velocities.
\end{enumerate}

% ============================================================
\section{Conclusion}
\label{sec:conclusion}
% ============================================================

This work demonstrated successful identification and control of the Heli2DoF platform. Empirical SISO models were identified from open-loop experiments and combined into a MIMO state-space representation. Both pole placement and LQR integral state-feedback regulators were designed and validated. The LQR REI controller met the transient and steady-state specifications in simulation, and achieved a settling time of 2.920\,s and an overshoot of 2.4\% experimentally, satisfying all critical performance requirements. Future work should address the observed steady-state error under PASCAL coordinate transformation and improve yaw disturbance suppression.

% ============================================================
% REFERENCES
% ============================================================
\begin{thebibliography}{9}

\bibitem{labcontrol2020}
E. Interiano,
``Proyecto Corto I-2020: Control de motor CD,''
Lab.\ de Control Automático, ITCR, 2020.
[Online]. Available: \url{http://www.ie.tec.ac.cr/einteriano/control/Laboratorio/4.2ProyectoCorto_I_2020__motorCD.pdf}

\bibitem{video}
E. Interiano,
``Experimento Heli2DoF,'' YouTube, 2020.
[Online]. Available: \url{https://youtu.be/3yfB-GlOzhk}

\bibitem{heli_zip}
E. Interiano,
``Heli2DoF data and LabControl application,'' ITCR, 2020.
[Online]. Available: \url{http://www.ie.tec.ac.cr/einteriano/control/Laboratorio/4.5Heli2DoF.zip}

\bibitem{pascal_zip}
E. Interiano,
``PASCAL application,'' ITCR, 2020.
[Online]. Available: \url{http://www.ie.tec.ac.cr/einteriano/control/Laboratorio/4.2AppLabControl.zip}

\bibitem{ogata}
K. Ogata,
\textit{Modern Control Engineering}, 5th ed.
Prentice Hall, 2010.

\bibitem{franklin}
G. F. Franklin, J. D. Powell, and A. Emami-Naeini,
\textit{Feedback Control of Dynamic Systems}, 8th ed.
Pearson, 2019.

\end{thebibliography}

\end{document}